\documentclass[11pt]{amsart}
\usepackage[T1]{fontenc}
\usepackage[utf8]{inputenc}
\usepackage{amssymb}
\usepackage[margin=1.1in]{geometry}
\usepackage[hidelinks]{hyperref}
\newcommand{\stamptime}{\the\numexpr\time/60\relax:\ifnum\numexpr\time-60*(\time/60)\relax<10 0\fi\the\numexpr\time-60*(\time/60)\relax}
\AddToHook{shipout/foreground}{\put(0.5\paperwidth,-\paperheight+0.55cm){\makebox[0pt]{\tiny\ttfamily unit/doc/\jobname.tex \ --- \ compiled \today, \stamptime}}}
\newtheorem{theorem}{Theorem}
\newtheorem{proposition}[theorem]{Proposition}
\newtheorem{lemma}[theorem]{Lemma}
\newtheorem{corollary}[theorem]{Corollary}
\theoremstyle{definition}
\newtheorem{definition}[theorem]{Definition}
\theoremstyle{remark}
\newtheorem{remark}[theorem]{Remark}
\newcommand{\Om}{\Omega}
\newcommand{\refl}{\mathsf{refl}}
\newcommand{\ap}{\mathsf{ap}}
\newcommand{\Fib}{\mathrm{Fib}}
\newcommand{\SCt}{\mathsf{SC}_3}
\newcommand{\SCf}{\mathsf{SC}_4}
\newcommand{\Tri}{\mathsf{Triples}}
\newcommand{\Mor}{\mathsf{Mor}}
\newcommand{\Hom}{\mathrm{Hom}}
\newcommand{\Map}{\mathrm{Map}_*}
\newcommand{\real}{\mathsf{realise}}
\newcommand{\rmap}{\rho}
\newcommand{\lp}{\mathsf{loop}}
\newcommand{\Unloop}{\mathsf{Unloop}}
\newcommand{\Tl}{\mathsf{Tr}}
\newcommand{\Z}{\mathbb{Z}}
\newcommand{\cls}[1]{[#1]}
\newcommand{\cmp}{\mathbin{;}}
\title{Arbitrary pointed maps between Postnikov triples,\\ and the spheres}
\date{28 October 2026}
\begin{document}
\maketitle
\setlength{\emergencystretch}{2em}

\begin{abstract}
In homotopy type theory, pointed simply connected $3$-types are classified by
triples $(A,B,k)$ with $k:K(A,2)\to_*K(B,4)$. We show that \emph{arbitrary}
pointed maps between their realisations are classified as well: the canonical
map from triples of a homomorphism $a$, a homomorphism $b$ and a pointed
homotopy $K(a,2)\cmp k'=k\cmp K(b,4)$ to pointed maps of fibres is an
equivalence of types. It follows that the homotopy precategory of triples is
isomorphic to that of pointed simply connected $3$-types. We then compute the
universal properties of $S^2$ and, through the looping embedding of pointed
$2$-connected $4$-types, of $S^3$ on the side of triples, and recover
$\pi_3(S^2)\cong\Z$ and $\pi_4(S^3)\cong\Z/2$. All results are checked
formally; this version omits the correspondence with the formal declarations,
which is in the companion note.
\end{abstract}

\noindent\textbf{Attribution.} The strategy (classify objects and arbitrary
maps, identify the sphere universal objects on the side of triples, and treat
$\pi_4(S^3)$ through the looping embedding) and the theory of deloopings used
below are David Wärn's. The classification of objects is Hart--Milner's (Owen
Milner's equivalence). We do not know Owen Milner's unpublished treatment of
maps and do not claim to reproduce it.

\noindent\textbf{Conventions.} A type is \emph{$n$-connected} or an
\emph{$n$-type} in the usual (homotopical) indexing. Composition is
diagrammatic: $f\cmp g$ is ``$f$, then $g$''; for pointed maps
$(f\cmp g)_1=g_1\circ f_1$ and $(f\cmp g)_0=\ap_{g_1}(f_0)\cdot g_0$. Path
composition $p\cdot q$ is defined by induction on $q$, so $p\cdot\refl=p$
holds by definition. $K(A,n)$ denotes the Eilenberg--MacLane space and
$K(h,n)$ the map induced by a homomorphism $h$.

\section{The statement}

A \emph{triple} is $t=(A,B,k)$ with abelian groups $A,B$ and a pointed map
$k:K(A,2)\to_*K(B,4)$. Its realisation is the fibre
\[\Fib k=\textstyle\sum_{x:K(A,2)}(k\,x=\mathrm{pt}),\]
pointed by $(\mathrm{pt},k_0)$; it is a pointed simply connected $3$-type, and
$\real:\Tri\to\SCt$ is an equivalence of types (Hart--Milner). Fix triples
$t=(A,B,k)$, $t'=(A',B',k')$ and write $V=K(A,2)$, $Y=K(B,4)$, $V'=K(A',2)$,
$Y'=K(B',4)$, $P_x=(k_1x=y_0)$, and $p:\Fib k\to_*V$ for the projection. We use:
$V$ is simply connected and a $2$-type, $Y$ and $Y'$ are $3$-connected
$4$-types, so each $P_x$ is $2$-connected and each path type of $Y'$ is a
$3$-type.

\begin{definition}\label{def:mor}
\[\Mor(t,t')=\sum_{a:\Hom(A,A')}\ \sum_{b:\Hom(B,B')}
 \big(K(a,2)\cmp k'=_{\Map(V,Y')}k\cmp K(b,4)\big).\]
The third component is a \emph{path of pointed maps}: a homotopy together with
its coherence at the base point.
\end{definition}

\begin{definition}\label{def:rmap}
For $(a,b,H)$ put $u=K(a,2)$, $v=K(b,4)$, let $hs_x:k'_1(u_1x)=v_1(k_1x)$ be the
underlying homotopy of $H$ and $c:\ap_{k'_1}(u_0)\cdot k'_0=hs_{v_0}\cdot(\ap_{v_1}k_0\cdot v_0)$
its coherence. The \emph{canonical map}
$\rmap(a,b,H):\Fib k\to_*\Fib k'$ is
\[(x,q)\longmapsto\big(u_1x,\ hs_x\cdot(\ap_{v_1}q\cdot v_0)\big),\]
pointed by the pair of $u_0$ and the path obtained from $c$.
\end{definition}

\begin{theorem}\label{thm:main}
For all triples $t,t'$, the canonical map $\rmap:\Mor(t,t')\to\Map(\Fib k,\Fib k')$
is an equivalence of types.
\end{theorem}

\noindent No hypothesis on the groups or on $k,k'$ is needed.

\begin{remark}
(1) Both sides are $1$-types and, in general, not sets: $\Map(V,Y')$ is a
$2$-type, so its path types are $1$-types; $\Map(\Fib k,\Fib k')$ is a
$1$-type since $\Fib k$ is simply connected and $\Fib k'$ is a $3$-type.

(2) The homotopy is data. For $t=(A,0,0)$, $t'=(0,C,0)$ the coefficients are
trivial and both sides are $\Map(K(A,2),K(C,3))$, whose set of components is
$H^3(K(A,2);C)$, non-zero in general.

(3) Classically the theorem follows from a connectivity estimate for the
cofibre of $\Fib k\to K(A,2)$. We use no cofibres; the key lemma below is a
direct fibrewise replacement of that estimate.
\end{remark}

\section{Proof of Theorem~\ref{thm:main}}

For a pointed $c=(c_1,c_0):V\to_*Y'$ put $Q^c_x=(c_1x=y'_0)$ and
\[T(c)=\sum_{F:\prod_x P_x\to Q^c_x}\big(F_{v_0}k_0=c_0\big),\qquad
S(c)=\sum_{\beta:\Map(Y,Y')}\big(c=_{\Map(V,Y')}k\cmp\beta\big).\]
$T(c)$ is the type of pointed null-homotopies of $p\cmp c$ in fibrewise form,
and $S(c)$ the type of factorisations of $c$ through $k$.

\subsection{Maps out of $\Fib k$ into $\Fib k'$}

\begin{lemma}\label{lem:dec}
The map $\mathrm{Dec}:\sum_{\alpha:\Map(V,V')}T(\alpha\cmp k')\to\Map(\Fib k,\Fib k')$,
\[(\alpha,F,e)\longmapsto\big((x,q)\mapsto(\alpha_1x,F_xq),\ (\alpha_0,\pi(e))\big),\]
is an equivalence. Here $e:F_{v_0}k_0=\ap_{k'_1}\alpha_0\cdot k'_0$ (which is
$(\alpha\cmp k')_0$ by definition), and $\pi(e)$ is the path
$\alpha_0{}_*(F_{v_0}k_0)=(\ap_{k'_1}\alpha_0)^{-1}\cdot F_{v_0}k_0=(\ap_{k'_1}\alpha_0)^{-1}\cdot(\ap_{k'_1}\alpha_0\cdot k'_0)=k'_0$.
\end{lemma}
\begin{proof}
(i) Precomposition with $p$ is an equivalence $(V\to V')\simeq(\Fib k\to V')$:
the fibres of $p$ are the $P_x$, which are $2$-connected, and $V'$ is a
$2$-type.
(ii) Hence $(a_1,F)\mapsto((x,q)\mapsto(a_1x,F_xq))$ is an equivalence
\[\textstyle\sum_{a_1:V\to V'}\prod_x P_x\to(k'_1(a_1x)=y'_0)\ \simeq\ (\Fib k\to\Fib k'),\]
by currying, (i), and the distributivity of $\Pi$ over $\Sigma$.
(iii) A pointed map is a map with a path $G(w_0)=w'_0$; along (ii) this becomes
the path $(a_1v_0,F_{v_0}k_0)=w'_0$, by definition.
(iv) Paths in a $\Sigma$-type: for fixed $(a_1,F)$,
$(\alpha_0,e)\mapsto(\alpha_0,\pi(e))$ is an equivalence onto
$(a_1v_0,F_{v_0}k_0)=w'_0$, since $\pi$ composes $\ap$ of left multiplication by a
fixed path (an equivalence, so its $\ap$ is one) with a left and a right
multiplication by fixed paths.
Reordering the components gives $\mathrm{Dec}$.
\end{proof}

\subsection{The key lemma}

For $\beta:\Map(Y,Y')$ write $G^\beta_y(q)=\ap_{\beta_1}q\cdot\beta_0:\beta_1y=y'_0$
for $q:y=y_0$ (so $(k\cmp\beta)_0=G^\beta_{k_1v_0}(k_0)$), and
$Z^\beta_x=(c_1x=\beta_1(k_1x))$. Define
\[\Phi_c:S(c)\to T(c),\qquad \Phi_c(\beta,s)=\big(x\,q\mapsto hs_x\cdot G^\beta_{k_1x}(q),\ c(s)^{-1}\big),\]
where $hs$ is the homotopy of $s$ and $c(s):c_0=hs_{v_0}\cdot(k\cmp\beta)_0$ its
coherence.

\begin{lemma}[key lemma]\label{lem:key}
For every pointed $c:V\to_*Y'$, $\Phi_c$ is an equivalence: pointed
null-homotopies of $p\cmp c$ are the same as factorisations of $c$ through $k$.
\end{lemma}

For $F_x:P_x\to Q^c_x$ say that $F_x$ is a \emph{translate} of $G^\beta$ if we
are given
\[\Tl^\beta_x(F_x)=\sum_{z:Z^\beta_x}\big((q\mapsto z\cdot G^\beta_{k_1x}q)=F_x\big).\]

\begin{lemma}[paths of pointed maps]\label{lem:ptpath}
For pointed $c,g:X\to_*Y$, sending a path to its homotopy and coherence is an
equivalence $(c=g)\simeq\sum_{h:\prod_xc_1x=g_1x}(c_0=h_{x_0}\cdot g_0)$.
\end{lemma}
\begin{proof}
Fundamental theorem of identity types: the total space
$\sum_g\sum_h(c_0=h_{x_0}\cdot g_0)$ is contractible (first a singleton of
$c_1$ by function extensionality, then a singleton of $c_0$), and the map
agrees with the transport map by path induction.
\end{proof}

\begin{lemma}[translates]\label{lem:transl}
For each $x$, $z\mapsto(q\mapsto z\cdot G^\beta q,\ (z,\refl))$ is an
equivalence $Z^\beta_x\simeq\sum_{F_x}\Tl^\beta_x(F_x)$, and each
$\Tl^\beta_x(F_x)$ is a proposition.
\end{lemma}
\begin{proof}
The inverse forgets $F_x$ and the path; one round trip is trivial, the other is
path induction. $\Tl^\beta_x(F_x)$ is the fibre over $F_x$ of
$\tau:z\mapsto(q\mapsto z\cdot G^\beta q)$, which is the constant map
$Z^\beta_x\to(P_x\to Z^\beta_x)$ followed by postcomposition with the fibrewise
equivalence $d\mapsto(q\mapsto d\,q\cdot G^\beta q)$ (right multiplication). The
constant map is an embedding because $P_x$ is $2$-connected and $Z^\beta_x$ a
$3$-type (maps from a $2$-connected type into the path types, which are
$2$-types, are constant); so $\tau$ is an embedding and its fibres are
propositions.
\end{proof}

\begin{lemma}\label{lem:fac}
For fixed $\beta$, the map
$(c=k\cmp\beta)\to\sum_{(F,e):T(c)}\prod_x\Tl^\beta_x(F_x)$,
$s\mapsto(\Phi_c(\beta,s),\ x\mapsto(hs_x,\refl))$, is an equivalence.
\end{lemma}
\begin{proof}
Lemma~\ref{lem:ptpath} gives
$(c=k\cmp\beta)\simeq\sum_{h:\prod_xZ^\beta_x}(c_0=h_{v_0}\cdot G^\beta_{k_1v_0}k_0)$.
Lemma~\ref{lem:transl} under $\prod_x$ and distributivity give
$\prod_xZ^\beta_x\simeq\sum_F\prod_x\Tl^\beta_x(F_x)$, with
$F_{v_0}k_0=h_{v_0}\cdot G^\beta_{k_1v_0}k_0$ by definition; the pointing
condition is carried along, and inverted.
\end{proof}

The restriction of $(F,e):T(c)$ to the fibre over the base point, conjugated
by $c_0$, is the pointed map
\[\mathrm{res}(F,e):\Om Y\to_*\Om Y',\qquad q\mapsto c_0^{-1}\cdot F_{v_0}(k_0\cdot q),\]
pointed by $c_0^{-1}\cdot F_{v_0}k_0=c_0^{-1}\cdot c_0=\refl$.

\begin{lemma}\label{lem:atpt}
For $(F,e):T(c)$ and $\beta:\Map(Y,Y')$, the propositions
$\Tl^\beta_{v_0}(F_{v_0})$ and $\mathrm{res}(F,e)=\lp\,\beta$ are logically
equivalent.
\end{lemma}
\begin{proof}
The second is a proposition since $\Map(\Om Y,\Om Y')$ is a set ($\Om Y$ is
$2$-connected, $\Om Y'$ a $3$-type). We use
$\lp\,\beta(q)=\beta_0^{-1}\cdot(\ap_{\beta_1}q\cdot\beta_0)$.

($\Rightarrow$) Given $(z,E)$, $F_{v_0}w=z\cdot(\ap_{\beta_1}w\cdot\beta_0)$ for
all $w$, and with $e$, $c_0=z\cdot(\ap_{\beta_1}k_0\cdot\beta_0)$. Since
$\Om^2Y'$ is connected, pointed maps $\Om Y\to_*\Om Y'$ with equal underlying
maps are merely equal; as the goal is a proposition it suffices to check, for
all $q$,
\[(z\cdot(\ap_{\beta_1}k_0\cdot\beta_0))^{-1}\cdot\big(z\cdot(\ap_{\beta_1}(k_0\cdot q)\cdot\beta_0)\big)
=\beta_0^{-1}\cdot(\ap_{\beta_1}q\cdot\beta_0).\]
This follows from $(z\cdot A)^{-1}\cdot(z\cdot B)=A^{-1}\cdot B$,
$\ap_{\beta_1}(k_0\cdot q)=\ap_{\beta_1}k_0\cdot\ap_{\beta_1}q$, and
$(A\cdot b)^{-1}\cdot((A\cdot M)\cdot b)=b^{-1}\cdot(M\cdot b)$ (path induction on
$z$ and on $b$).

($\Leftarrow$) Given $r:\mathrm{res}(F,e)=\lp\,\beta$, take
$z=F_{v_0}k_0\cdot(\ap_{\beta_1}k_0\cdot\beta_0)^{-1}$. For $w=k_0\cdot q$ we get
$F_{v_0}w=c_0\cdot\mathrm{res}(F,e)(q)=c_0\cdot\lp\,\beta(q)$, and with $e$ the same
identities give $z\cdot(\ap_{\beta_1}w\cdot\beta_0)=F_{v_0}w$. Every $w$ is of
this form ($w=k_0\cdot(k_0^{-1}\cdot w)$); function extensionality gives $E$.
\end{proof}

\begin{lemma}\label{lem:unique}
For every $(F,e):T(c)$, the type $\sum_\beta\prod_x\Tl^\beta_x(F_x)$ is
contractible.
\end{lemma}
\begin{proof}
For each $\beta$, $\prod_x\Tl^\beta_x(F_x)$ and $\Tl^\beta_{v_0}(F_{v_0})$ are
logically equivalent propositions: one way by evaluation, the other because a
family of propositions over the connected $V$ holding at $v_0$ holds
everywhere. With Lemma~\ref{lem:atpt}, the type is equivalent to
$\sum_\beta(\mathrm{res}(F,e)=\lp\,\beta)$, the type of deloopings of
$\mathrm{res}(F,e)$. It is contractible by Wärn's delooping theorem: since $Y$ is
$2$-connected and $Y'$ a $4$-type, every pointed map $\Om Y\to_*\Om Y'$ has a
unique delooping.
\end{proof}

\begin{proof}[Proof of Lemma~\ref{lem:key}]
$\Phi_c$ is the composite
\[S(c)\simeq\sum_\beta\sum_{(F,e)}\prod_x\Tl^\beta_x(F_x)
\simeq\sum_{(F,e)}\sum_\beta\prod_x\Tl^\beta_x(F_x)\longrightarrow T(c)\]
of Lemma~\ref{lem:fac} under $\sum_\beta$, a swap, and the first projection,
which is an equivalence by Lemma~\ref{lem:unique}.
\end{proof}

\begin{remark}
A fibrewise map $F_x:P_x\to Q^c_x$ between two ``$K(B,3)$-torsors'' is a linear
part and a translation. The linear part is the delooping $\beta$ of the
restriction over $v_0$; the translation is the homotopy component, kept as
data and fixed at $v_0$ by the pointing. No point is chosen in any other fibre,
no representative of a class is chosen, and lifts are not assumed unique.
\end{remark}

\subsection{Coefficients and assembly}

\begin{lemma}\label{lem:coef}
The maps $a\mapsto K(a,2)$, $\Hom(A,A')\to\Map(V,V')$, and $b\mapsto K(b,4)$,
$\Hom(B,B')\to\Map(Y,Y')$, are equivalences.
\end{lemma}
\begin{proof}
At level $2$ this is the known classification of maps $K(A,2)\to_*K(A',2)$. At
level $4$, a map $G$ is $K(\mathsf{inv}(G),4)$ by climbing from level $2$ with
unique deloopings, and $\mathsf{inv}(K(b,4))=b$ because the looping of
$K(b,m+1)$ is the conjugate of $K(b,m)$ by the equivalences
$\Om K(-,m+1)\simeq K(-,m)$, and conjugation is invertible.
\end{proof}

\begin{proof}[Proof of Theorem~\ref{thm:main}]
$\rmap$ is the composite
\begin{align*}
\Mor(t,t')&\simeq\textstyle\sum_aS(K(a,2)\cmp k')
\xrightarrow{\ \sum_a\Phi\ }\textstyle\sum_aT(K(a,2)\cmp k')\\
&\simeq\textstyle\sum_{\alpha}T(\alpha\cmp k')
\xrightarrow{\ \mathrm{Dec}\ }\Map(\Fib k,\Fib k'),
\end{align*}
using Lemma~\ref{lem:coef} at the first and third steps and
Lemmas~\ref{lem:key} and~\ref{lem:dec}. That the composite is the canonical
map of Definition~\ref{def:rmap} holds by definition, including the pointing.
\end{proof}

\begin{remark}
The canonical map lies over $K(a,2)$, $\rmap(a,b,H)\cmp p'=p\cmp K(a,2)$, and
its restriction to the fibre over the base point is $\lp\,K(b,4)$
(Lemma~\ref{lem:atpt}). That $b$ is the map induced on $\pi_3$ under
$\pi_3(\Fib k)\cong B$ is not proved here.
\end{remark}

\section{The homotopy precategory of triples}\label{sec:cat}

\begin{definition}
For $m=(a,b,H):t\to t'$ and $m'=(a',b',H'):t'\to t''$ let
$m\cmp m'=(a\cmp a',\ b\cmp b',\ H'')$, with $H''$ the pasting
\[K(a\cmp a',2)\cmp k''=K(a,2)\cmp K(a',2)\cmp k''=K(a,2)\cmp k'\cmp K(b',4)
=k\cmp K(b,4)\cmp K(b',4)=k\cmp K(b\cmp b',4),\]
built from $K(g\cmp h,n)=K(g,n)\cmp K(h,n)$, associativity, $H'$ whiskered by
$K(a,2)$, and $H$ whiskered by $K(b',4)$. The identity is
$(1_A,1_B,H_1)$, where $H_1$ comes from $K(1,n)=\mathrm{id}$ and the unit laws.
\end{definition}

The identity $K(1,n)=\mathrm{id}$ ($n=2,4$) follows from surjectivity of
$K(-,n)$: if $K(h,n)=\mathrm{id}$, then
$K(1,n)=K(1,n)\cmp K(h,n)=K(1\cmp h,n)=K(h,n)=\mathrm{id}$.

\begin{proposition}\label{prop:funct}
$\rmap(\mathrm{id}_t)$ and the identity, and $\rmap(m\cmp m')$ and
$\rmap(m)\cmp\rmap(m')$, are pointwise equal; hence they are equal as homotopy
classes of pointed maps.
\end{proposition}
\begin{proof}
For a path $r$ of pointed maps let $hp(r)$ be its underlying homotopy. By path
induction, $hp$ takes composites to composites and inverses to inverses,
$hp(\ap_{u\cmp-}s)_x=hp(s)_{u_1x}$, $hp(\ap_{-\cmp v}s)_x=\ap_{v_1}(hp(s)_x)$,
the associativity path has trivial homotopy, and so has a path changing only the
pointing. Then: (1) the canonical map of a square transported along paths
$u_2=u$, $v_2=v$ is pointwise the canonical map of the square (path induction
on the two paths); (2) for a pasted square, the homotopy is
$hs'_{u_1x}\cdot\ap_{v'_1}(hs_x)$, and the second components agree by
$\ap$ of composites and reassociation; (3) the canonical map of the identity
square is pointwise the identity. Pointwise equality gives mere equality of
pointed maps, because the loop space of the target is connected.
\end{proof}

\noindent No coherent pointed path is constructed; the homotopy precategory only
needs classes.

\begin{theorem}\label{thm:hocat}
Let $\hom(t,t')=\|\Mor(t,t')\|_0$, with the composition and identities induced
by the definition above. This is a precategory, and the functor which is
$\real$ on objects and the class of $\rmap$ on hom-sets, into the homotopy
precategory of pointed simply connected $3$-types, is an \emph{isomorphism of
precategories}: an equivalence of object types and a bijection on hom-sets,
preserving identities and composition.
\end{theorem}
\begin{proof}
The class of $\rmap$ is a bijection by Theorem~\ref{thm:main}, and $\real$ is an
equivalence of types. By Proposition~\ref{prop:funct} it preserves composition
and identities, and the category laws in $\hom$ follow from those of pointed
maps up to homotopy through the bijections.
\end{proof}

\begin{remark}
Following Ahrens, Kapulkin and Shulman (\emph{Univalent categories and the Rezk
completion}, Definition~6.9), a precategory may have an arbitrary type of
objects, here $2$-types, while its hom-types are sets. Neither precategory is
univalent, and no such claim is made. By their Lemma~6.16, under univalence an
isomorphism of precategories yields an equality of precategories.
\end{remark}

\section{The components of the compatible homotopies}\label{sec:comp}

The sphere arguments compute $\hom(t_0,t')$ directly from $\Mor(t_0,t')$. The
coefficient $b$ is fixed by cohomological information about $K(\Z,2)$, and the
components of the type of homotopies by (H3). The latter concerns loops at an
arbitrary point of $\Map(K(\Z,2),K(B',4))$, which we move to the base point.

\begin{lemma}[translation]\label{lem:translation}
Let $X$, $Z$ be pointed, $M=\Map(X,\Om Z)$ and $G:M$. Pointwise right
multiplication by $G^{-1}$ is an equivalence $M\to M$ sending $G$ to a map merely
equal to the constant map, when $\Om^2Z$ is connected.
\end{lemma}

\begin{lemma}[components]\label{lem:components}
Let $Y\simeq_*\Om Z$ with $\Om^2Z$ connected, and suppose every pointed map
$X\to_*\Om^2Z$ is merely constant. Then any two paths $f=g$ in $\Map(X,Y)$ are
merely equal.
\end{lemma}
\begin{proof}
Transport along $Y\simeq_*\Om Z$ and translate by $g^{-1}$
(Lemma~\ref{lem:translation}); this reduces to two paths $r,r':u=v$ with $v$
merely, hence (the goal being a proposition) actually, the constant map. Then
$r^{-1}\cdot r'$ is a loop at the constant map, that is, a pointed map
$X\to_*\Om^2Z$, merely constant by hypothesis.
\end{proof}

\noindent For $X=K(\Z,2)$ and $Y=K(B',4)\simeq_*\Om K(B',5)$, the hypothesis is
(H3): $[K(\Z,2),K(B',3)]_*$ is a point, since $\Om^2K(B',5)\simeq_*K(B',3)$.

\section{$\pi_3(S^2)$ through triple morphisms}\label{sec:pi3}

Let $\kappa:K(\Z,2)\to_*K(\Z,4)$ be the cup square, $t_0=(\Z,\Z,\kappa)$ and
$P=\Fib\kappa$. We use:
\begin{itemize}
\item (H4): $b\mapsto\cls{\kappa\cmp K(b,4)}$ is a bijection
$\Hom(\Z,B')\to[K(\Z,2),K(B',4)]_*$;
\item (H3): $[K(\Z,2),K(B',3)]_*$ is a point.
\end{itemize}

\begin{proposition}\label{prop:theta}
For every triple $t'=(A',B',k')$, the map $\theta:\hom(t_0,t')\to\Hom(\Z,A')$,
$\cls{a,b,H}\mapsto a$, is a bijection.
\end{proposition}
\begin{proof}
\emph{Surjective:} for $a$, (H4) gives $b$ with
$\cls{\kappa\cmp K(b,4)}=\cls{K(a,2)\cmp k'}$, hence merely a homotopy $H$.
\emph{Injective:} if two morphisms lie over $a$, their coefficients $b,b'$ have
the same image under (H4), so $b=b'$, and their homotopies are merely equal by
Lemma~\ref{lem:components}.
\end{proof}

\begin{proposition}\label{prop:mark3}
Let $s\in[S^2,P]_*$ be the class with $s\cmp\cls p=\cls\sigma$, where
$\sigma:S^2\to_*K(\Z,2)$ is the generator. For every triple $t'$,
$s^*:[P,\Fib k']_*\to[S^2,\Fib k']_*$ is a bijection.
\end{proposition}
\begin{proof}
The composite
\[\hom(t_0,t')\xrightarrow{\cls\rmap}[P,\Fib k']_*\xrightarrow{s^*}[S^2,\Fib k']_*\xrightarrow{p'_*}[S^2,K(A',2)]_*\]
sends $\cls m$ to $\cls{\sigma\cmp K(a,2)}$, because $\rmap(m)$ lies over
$K(a,2)$. The map $a\mapsto\cls{\sigma\cmp K(a,2)}$ is a bijection
$\Hom(\Z,A')\to[S^2,K(A',2)]_*$: maps $S^2\to K(A',2)$ correspond to maps
$K(\Z,2)\to K(A',2)$ through $\sigma$, and these to homomorphisms. So the
composite is a bijection. The first map is a bijection by
Theorem~\ref{thm:main}, and $p'_*$ by lifting from the source $S^2$ (maps from
$S^2$ into $\Fib k'$ correspond to maps into $K(A',2)$). Hence $s^*$ is a
bijection.
\end{proof}

\begin{corollary}
$\pi_3(S^2)\cong\Z$, from the universal property of $S^2$.
\end{corollary}
\begin{proof}
By the classification of objects, every pointed simply connected $3$-type is
$\Fib k'$ for some triple, so $s^*$ is a bijection for all of them. Then $P$ is
the $3$-truncation of $S^2$ (a two-map Yoneda argument), and
$\pi_3(S^2)\cong\pi_3(P)\cong\Z$.
\end{proof}

\section{$\pi_4(S^3)$ through triple morphisms and the looping embedding}\label{sec:pi4}

Looping is an embedding of pointed $2$-connected $4$-types into pointed simply
connected $3$-types; its image is the \emph{deloopable} triples. On maps,
looping is a bijection $[X,Y]_*\simeq[\Om X,\Om Y]_*$ for $X$ $2$-connected and
$Y$ a $4$-type. Let $\bar k=\kappa\cmp K(\rho,4):K(\Z,2)\to_*K(\Z/2,4)$ be the
cup square reduced mod~$2$, $\bar t=(\Z,\Z/2,\bar k)$, $R=\Fib\bar k$, and $BR$
its delooping. The universal property to establish is statement M: for a
marked class $s\in[S^3,BR]_*$, the map $s^*:[BR,Y]_*\to[S^3,Y]_*$ is a bijection
for every pointed $2$-connected $4$-type $Y$. By the bijection on maps, M at
$Y$ is equivalent to bijectivity of $r^*:[R,\Om Y]_*\to[\Om S^3,\Om Y]_*$ for the
looped class $r$.

\begin{proposition}\label{prop:theta4}
For every deloopable triple $t'=(A',B',k')$, the map
$\theta:\hom(\bar t,t')\to\Hom(\Z,A')$ is a bijection.
\end{proposition}
\begin{proof}
\emph{Surjective:} by (H4), $\cls{K(a,2)\cmp k'}=\cls{\kappa\cmp K(h,4)}$ for
some $h$. Since $t'$ is deloopable, $2h(1)=0$ (Statement~T, the restricted
$2$-torsion lemma). So $h=\rho\cmp g$, and
$\cls{K(a,2)\cmp k'}=\cls{\bar k\cmp K(g,4)}$.
\emph{Injective:} $g\mapsto\cls{\bar k\cmp K(g,4)}$ is injective, which fixes
$b$; the homotopies are merely equal by Lemma~\ref{lem:components} and (H3).
\end{proof}

\begin{proposition}
Statement M holds for every pointed $2$-connected $4$-type $Y$.
\end{proposition}
\begin{proof}
Let $t'$ be the triple of $\Om Y$; it is deloopable.

\emph{Lifting from $\Om S^3$.} A map $F:\Om S^3\to_*K(A',2)$ is, up to homotopy,
$r\cmp\bar p\cmp K(a,2)$ for some $a$. This uses that $r^*$ is bijective into
$K(A',2)$ (M at $K(A',3)$, looped), that $\bar p^*$ is bijective on maps into
$2$-types, and Lemma~\ref{lem:coef}. By Proposition~\ref{prop:theta4} there is
merely a morphism $(a,b,H)$, so $F\cmp k'$ is merely
$r\cmp\bar p\cmp\bar k\cmp K(b,4)$, which is merely constant. Together with
uniqueness of lifts from $\Om S^3$, maps $\Om S^3\to\Fib k'$ correspond to maps
$\Om S^3\to K(A',2)$.

\emph{The square.} The composite
\[\hom(\bar t,t')\to[R,\Fib k']_*\xrightarrow{r^*}[\Om S^3,\Fib k']_*\to[\Om S^3,K(A',2)]_*\]
sends $\cls m$ to $r\cmp\cls{\bar p\cmp K(a,2)}$, and
$a\mapsto r\cmp\cls{\bar p\cmp K(a,2)}$ is a bijection. So $r^*$ is bijective at
$\Fib k'\simeq\Om Y$, and M holds at $Y$.
\end{proof}

\begin{corollary}
$\pi_4(S^3)\cong\Z/2$, from the universal property of $S^3$.
\end{corollary}
\begin{proof}
Apply statement M at $Y=BR$ and at $Y=\|S^3\|_4$; the two-map Yoneda argument
gives $\pi_4(S^3)\cong\pi_4(BR)\cong\pi_3(R)\cong\Z/2$.
\end{proof}

\section{What is used}

The sphere arguments rest on Theorem~\ref{thm:main} and the computation of
hom-sets from $\Mor$. They reuse the following:
\begin{itemize}
\item the cup square and its normalisation;
\item (H3) and (H4) for $K(\Z,2)$;
\item the restricted $2$-torsion lemma;
\item Wärn's delooping theorem and the looping embedding;
\item the elementary computations of maps out of $S^2$, $S^3$ and $\Om S^3$,
together with the marked classes.
\end{itemize}
They do not use any earlier, specialised mapping theorem for the targets
$\Fib k$: this was confirmed formally by deleting those results from a copy of
the development, which still checks.

\end{document}
